\documentclass[11pt]{article}
\usepackage[margin=1in]{geometry}
\usepackage{amsmath,amssymb,amsthm}
\usepackage{array,booktabs,graphicx,tabularx}
\usepackage{fancyvrb}
\usepackage{hyperref}
\usepackage{xurl}
\hypersetup{colorlinks=true,linkcolor=black,urlcolor=blue}
\setlength{\parindent}{0pt}
\setlength{\parskip}{0.55em}
\title{An All-Obtuse Convex Hexagon Compiler for Planar Shifts of Finite Type\\[0.4em]\large Exact decoding under unrestricted contacts}
\author{}
\date{Referee draft, 27 July 2026}
\begin{document}
\maketitle


\begin{abstract}

We give an effective construction that maps every nearest-neighbor
\(\mathbb Z^2\)-shift of finite type \(S\) to a finite protoset \(P(S)\) of
unmarked strictly convex rational hexagons. Every interior angle of every
output hexagon lies in \((\pi/2,\pi)\). Tiles may be translated, rotated, or
reflected, and contacts are unrestricted. Even under this permissive model,
every \(P(S)\)-tiling is forced to be edge-to-edge, its abstract cell complex
is the standard honeycomb, and it locally decodes to a point of \(S\).
Conversely, every point of \(S\) deforms the affine honeycomb into a
\(P(S)\)-tiling. Hence \(P(S)\) tiles the plane if and only if \(S\) is
nonempty.

The construction yields undecidability of tileability for finite sets of
strictly convex rational hexagons whose angles are all obtuse, with all
Euclidean isometries and arbitrary contacts allowed. If \(S\) is strongly
aperiodic, then \(P(S)\) is aperiodic in the same unrestricted shape-only
sense. As an explicit instance, the phase-two compilation of the 13
Kari-Culik Wang tiles has 4,956 pairwise noncongruent prototiles. An exact
finite artifact enumerates its states, graph edges, and coronas and checks
the finite inequalities used in the construction.

\end{abstract}
\section*{1. Introduction}

Matching rules are easy to state when tiles carry colors, arrows, or signed
edges. They are harder to enforce when the final objects are unmarked
polygons, reflections are permitted, and one cannot assume an edge-to-edge
tiling. A geometric encoding must first control arbitrary contacts, then
recover the intended local state at every vertex, and finally prove that no
other global assembly remains.

This paper gives such an encoding for any nearest-neighbor
\(\mathbb Z^2\)-shift of finite type. The output polygons have a deliberately
restricted form:

\begin{enumerate}
\item every tile is a strictly convex hexagon;
\item every coordinate is rational;
\item every interior angle is strictly obtuse;
\item every Euclidean isometry, including reflection, is permitted; and
\item no contact rule is imposed beyond disjoint interiors and coverage of the
   plane.
\end{enumerate}

The all-obtuse condition removes T-junctions. A finite phase carrier fixes the
three honeycomb edge directions. Locally legal source triples are collared at
honeycomb vertices. A mixed-radix horizontal displacement assigns a different
length to every compatible pair of collared endpoint states. At an arbitrary
physical vertex, the three decoded graph edges form a triangle in a line
graph. Bipartiteness forces that triangle to be a star, recovering one
collared state rather than three pairwise compatible fragments.

The main result is a compiler theorem, rather than a theorem restricted to
aperiodic inputs. Aperiodicity and undecidability follow as separate
corollaries.

\subsection*{Theorem 1.1. All-obtuse compiler}

For every nearest-neighbor \(\mathbb Z^2\)-shift of finite type \(S\), one can
effectively construct a finite set \(P(S)\) of unmarked strictly convex
rational hexagons such that:

\begin{enumerate}
\item every angle of every polygon in \(P(S)\) lies in \((\pi/2,\pi)\);
\item \(P(S)\) tiles \(\mathbb R^2\) if and only if \(S\neq\varnothing\);
\item every unrestricted \(P(S)\)-tiling is edge-to-edge and trivalent, with
   abstract cell complex the standard honeycomb;
\item its decoded phase map is the universal cover of \(H_2\) and its face
   symbols form a configuration in \(S\);
\item up to one global orthogonal map and translation, that tiling is exactly
   the bounded piecewise-affine deformation associated with its decoded
   source point; and
\item every configuration in \(S\) gives such an unrestricted
   \(P(S)\)-tiling.
\end{enumerate}

Here unrestricted means that translations, rotations, and reflections are
allowed and no edge-to-edge convention, selected atlas, matching relation, or
one-sidedness assumption is supplied.

\subsection*{Corollary 1.2. Undecidability}

It is undecidable whether a finite set of unmarked strictly convex rational
hexagons, all of whose angles lie in \((\pi/2,\pi)\), tiles the plane when all
Euclidean isometries and unrestricted contacts are allowed.

\subsection*{Corollary 1.3. Strong aperiodicity}

If \(S\) is nonempty and every point of \(S\) has no nonzero
\(\mathbb Z^2\)-period, then \(P(S)\) is aperiodic: it tiles the plane and no
admitted \(P(S)\)-tiling has a nonzero translational symmetry.

The priority boundary is narrow. Robinson announced a reduction to finite
sets of convex lattice-coordinate pentagons or hexagons in 1971, including
rotations and reflections, but the cited paper does not give that
construction or its proof. The contribution here is an explicit compiler
with the all-obtuse output restriction and a classification of every
unrestricted-contact tiling. We make no claim to have originated the broad
idea that convex polygons can carry undecidable tiling systems.

\section*{2. Source systems and polygon tilings}

\subsection*{2.1. Nearest-neighbor shifts}

Let \(A\) be a finite alphabet and let
\[
  R_E,R_N\subseteq A\times A
\]
be directed east and north adjacency relations. Define
\[
S=S(A,R_E,R_N)
\]
to be the set of all \(x\in A^{\mathbb Z^2}\) satisfying
\[
  (x_{i,j},x_{i+1,j})\in R_E,\qquad
  (x_{i,j},x_{i,j+1})\in R_N
\]
for every \((i,j)\in\mathbb Z^2\).

We place these symbols on the faces of a honeycomb. Two of the three
face-neighbor directions carry \(R_E\) and \(R_N\). The third direction is
unconstrained. This is an embedding of a square nearest-neighbor shift into
the triangular face lattice, not a claim about arbitrary three-relation
triangular-lattice shifts.

A point \(x\in S\) is periodic if some nonzero
\((u,v)\in\mathbb Z^2\) satisfies
\[
  x_{i+u,j+v}=x_{i,j}
\]
for all \(i,j\). We call \(S\) strongly aperiodic when it is nonempty and no
point of \(S\) is periodic in this sense.

\subsection*{2.2. Shape-only tilings}

A prototile is an unmarked simple polygon considered up to Euclidean
congruence. A \(P\)-tiling is a covering of \(\mathbb R^2\) by congruent
copies of members of a finite protoset \(P\), with pairwise disjoint
interiors. Reflections are allowed. Contacts may be partial, and an
edge-to-edge convention is not assumed. A finite protoset is allowed to be
empty; the empty protoset is a valid no-instance and does not tile.

This definition is the only placement rule used below. Symbols, phase
coordinates, and graph states are construction and proof devices. They are
not retained as marks on the polygons.

\section*{3. The finite phase carrier}

\subsection*{3.1. A four-face honeycomb torus}

Let \(H_2\) have vertices
\[
  B_{ij},W_{ij},\qquad i,j\in\mathbb Z/2\mathbb Z.
\]
The three neighbors of \(B_{ij}\) are
\[
  W_{ij},\qquad W_{i-1,j},\qquad W_{i,j-1}.
\]
Its four faces are the simple cycles
\[
\begin{split}
C_{ij}={}&(B_{ij},W_{ij},B_{i+1,j},W_{i+1,j-1},\\
          &B_{i+1,j-1},W_{i,j-1}).
\end{split}
\]
Thus \(H_2\) has 8 vertices, 12 edges, and 4 faces. It is a torus
cellulation, and its universal cover is the standard honeycomb.

Here is a direct check. Modulo two, the three displayed neighbors of every
black vertex are distinct. Each \(C_{ij}\) is a simple six-cycle; every edge
occurs in exactly two such cycles, with opposite induced orientations; and
the incident face sectors form one cyclic link at each vertex. The resulting
complex is therefore a connected closed orientable surface. Its Euler
characteristic is
\[
  8-12+4=0,
\]
so it is a torus. If the indices \(i,j\) are instead allowed to range over
\(\mathbb Z\), the same face cycles embed as the affine honeycomb of Section
4.1. Quotienting that planar complex by the translations
\((i,j)\mapsto(i+2,j)\) and \((i,j)\mapsto(i,j+2)\) gives \(H_2\). Hence the
planar honeycomb is its universal cover.

The modulus two phase has two jobs. It makes every face boundary a simple
six-cycle, and it ensures that the endpoints of a state-graph edge determine
one fixed phase edge rather than several parallel witnesses.

\subsection*{3.2. Collared vertex states}

Place a source face \(x_{ij}\) at every \((i,j)\in\mathbb Z^2\). At a black
honeycomb vertex use the ordered triple
\[
  B_{ij}(x_{ij},x_{i-1,j+1},x_{i-1,j}),
\]
and at a white vertex use
\[
  W_{ij}(x_{ij},x_{i,j+1},x_{i-1,j+1}).
\]

A black triple \(B_{ij}(a,b,c)\) is locally legal when
\[
  (c,a)\in R_E,\qquad (c,b)\in R_N.
\]
A white triple \(W_{ij}(a,d,b)\) is locally legal when
\[
  (a,d)\in R_N,\qquad (b,d)\in R_E.
\]
A vertex state is a phase vertex of \(H_2\) together with one such locally
legal ordered triple.

Order all black states and all white states effectively:
\[
  Q_B=\{q^B_0,\ldots,q^B_{B-1}\},\qquad
  Q_W=\{q^W_0,\ldots,q^W_{W-1}\}.
\]
If either set is empty, the construction returns the empty protoset. This
case cannot occur when \(S\) is nonempty.

\subsection*{3.3. Legal coronas}

Name a central source face \(a\) and its six cyclic neighbors
\[
  b,d,e,f,g,c.
\]
The seven symbols form a legal corona when
\[
  (c,a),(a,e),(b,d),(g,f)\in R_E
\]
and
\[
  (g,a),(a,d),(c,b),(f,e)\in R_N.
\]

For phase \(C_{ij}\), associate the six states
\[
\begin{split}
q_0&=B_{ij}(a,b,c),\\
q_1&=W_{ij}(a,d,b),\\
q_2&=B_{i+1,j}(e,d,a),\\
q_3&=W_{i+1,j-1}(f,e,a),\\
q_4&=B_{i+1,j-1}(f,a,g),\\
q_5&=W_{i,j-1}(g,a,c).
\end{split}
\]

Let \(G\) be the simple bipartite graph with vertex set
\[
  Q_B\sqcup Q_W.
\]
An unordered pair \(\{q,r\}\) is an edge of \(G\) when \(q,r\) occur
consecutively in at least one legal corona cycle. Multiple witnesses create
one graph edge.

Every state already contains its phase vertex. Hence a graph edge has a
fixed black phase endpoint, a fixed white phase endpoint, and one fixed
black-to-white phase vector. The graph has no loops or parallel edges.

\section*{4. Exact polygon geometry}

\subsection*{4.1. The affine base honeycomb}

Use the centrally symmetric hexagon
\[
\begin{aligned}
v_0&=(2,0),&
v_1&=(1,3/2),&
v_2&=(-1,3/2),\\
v_3&=(-2,0),&
v_4&=(-1,-3/2),&
v_5&=(1,-3/2).
\end{aligned}
\]
Its face-center lattice has basis
\[
  a=(-3,3/2),\qquad b=(0,3).
\]
The face at \((i,j)\) is the translate of the base face by
\[
  ia+jb.
\]
These translates form an affine image of the regular honeycomb.

\subsection*{4.2. Mixed-radix displacements}

Put
\[
  D=4096BW.
\]
Assign a horizontal displacement to every state:
\[
z(q)=
\begin{cases}
 b/D,&q=q^B_b,\\
 B(w+1)/D,&q=q^W_w.
\end{cases}
\]
For a legal corona cycle \((q_0,\ldots,q_5)\), define the polygon
\[
  P(q_0,\ldots,q_5)
  =
  \operatorname{poly}
  \bigl(v_0+(z(q_0),0),\ldots,v_5+(z(q_5),0)\bigr).
\]
Finally, \(P(S)\) is the finite set of congruence classes of these polygons
over all four phases and all legal coronas.

Every coordinate belongs to \((2D)^{-1}\mathbb Z\). Multiplication by
\(2D\) gives an integer-coordinate representative. The construction is
finite and effective from \(A,R_E,R_N\).

\subsection*{4.3. Every graph edge has a unique length}

Let \(q^B_bq^W_w\) be an edge of \(G\), and put
\[
  \delta=z(q^W_w)-z(q^B_b)
  =\frac{B(w+1)-b}{D}.
\]
Its numerator belongs to the block
\[
  \{Bw+1,\ldots,Bw+B\}.
\]
These blocks are pairwise disjoint, so \(\delta\) determines both \(b\) and
\(w\). Moreover,
\[
  0<\delta\leq\eta,\qquad \eta=\frac1{4096}.
\]

The squared lengths are
\[
  L_h(\delta)=4+4\delta+\delta^2
\]
for a horizontal phase edge and
\[
  L_s(\delta)=\frac{13}{4}-2\delta+\delta^2
\]
for either slant direction. The function \(L_h\) is strictly increasing on
\((0,\eta]\), and \(L_s\) is strictly decreasing there. Their ranges are
disjoint because
\[
  L_s(\delta)<\frac{13}{4}<4<L_h(\delta).
\]

A state pair fixes one edge of the simple phase graph \(H_2\), including the
sign of the slant component of its black-to-white vector. It cannot occur in
both slant directions. It follows that distinct graph edges have distinct
lengths.

\subsection*{Lemma 4.1. Corona cycles are pairwise noncongruent}

No two distinct legal corona cycles define congruent polygons.

\paragraph{Proof}

A congruence preserves the cyclic word of side lengths up to rotation and
reversal. Every length identifies one graph edge. Consecutive distinct edges
in a simple graph recover their unique common endpoint. The cyclic edge word
therefore reconstructs the cyclic state word. Two congruent polygons came
from the same corona cycle up to the dihedral presentation already absorbed
by congruence. \(\square\)

\subsection*{4.4. Quantitative convexity and obtuseness}

The following calculation supplies the deformation margin used later.

\subsection*{Lemma 4.2. Uniform all-obtuse bound}

For any legal corona and every assignment produced above, the polygon is
strictly convex and every interior angle is obtuse. More precisely, every
cross product of consecutive forward edges is at least
\[
  3-3\eta>0,
\]
and every dot product of consecutive forward edges is at least
\[
  \frac54>0.
\]
Equivalently, the dot product of the two rays from a polygon vertex to its
neighbors is at most \(-5/4\).

\paragraph{Proof}

Write \(\delta_k\in(0,\eta]\) for the black-to-white displacement difference
on side \(k\), with the index attached to the corresponding graph edge. In
cyclic order the six forward edge vectors have the forms
\[
\begin{split}
e_0&=(-1+\delta_0,3/2),\\
e_1&=(-2-\delta_1,0),\\
e_2&=(-1+\delta_2,-3/2),\\
e_3&=(1-\delta_3,-3/2),\\
e_4&=(2+\delta_4,0),\\
e_5&=(1-\delta_5,3/2).
\end{split}
\]
Their consecutive cross products are of two types:
\[
  \frac32(2+\delta_k)
\]
and
\[
  3-\frac32(\delta_k+\delta_{k+1}).
\]
They are at least \(3-3\eta\).

Their consecutive forward dot products are of two types:
\[
  (2+\delta_k)(1-\delta_{k+1})
\]
and
\[
  \frac94-(1-\delta_k)(1-\delta_{k+1}).
\]
The first is at least \(2(1-\eta)>5/4\), and the second is at least
\(5/4\). Positive cross products give strict convexity. At a vertex the two
interior rays are \(-e_{k-1}\) and \(e_k\), whose dot product is the negative
of the displayed forward dot product. Thus every interior angle lies in
\((\pi/2,\pi)\). \(\square\)

The same calculation applies to the three rays at a honeycomb vertex. In
their cyclic order, their pairwise cross products have the two forms just
displayed. Their cyclic order therefore survives every permitted
displacement.

\section*{5. Arbitrary contacts become a honeycomb}

\subsection*{Lemma 5.1. All admitted tilings are edge-to-edge and trivalent}

Let \(P\) be any finite set of strictly convex polygons whose interior angles
all lie in \((\pi/2,\pi)\). Every unrestricted \(P\)-tiling is edge-to-edge,
and exactly three tile corners meet at every tiling vertex.

\paragraph{Proof}

First note that the tiling is locally finite. Let \(r>0\) be the minimum
inradius and \(M<\infty\) the maximum diameter among the finitely many
prototiles. Every tile meeting a compact set \(K\) has a disjoint
radius-\(r\) inball contained in the bounded \(M\)-neighborhood of \(K\).
Area packing permits only finitely many such inballs.

Consider a point \(x\) in the relative interior of a side of one tile. That
tile contributes a sector of angle \(\pi\). All sectors contributed by tiles
on the other side must sum to \(\pi\).

If a tile corner occurred at \(x\), one corner would leave a positive angle
smaller than \(\pi/2\), which cannot be filled by another permitted corner
or by an edge-interior sector of angle \(\pi\). Two corners would already
sum to more than \(\pi\). Hence no tile corner can lie in the relative
interior of another tile side. The remaining half-plane is occupied by
exactly one edge-interior sector.

A change from one neighboring side to another along a side would create a
tile corner in its relative interior. Thus every side is paired along its
full length with one whole side.

At a tiling vertex, only polygon corners occur. Two corners sum to less than
\(2\pi\), while four corners sum to more than \(2\pi\). Exactly three
corners meet. \(\square\)

\subsection*{Lemma 5.2. Honeycomb regularization}

The abstract cell complex of any unrestricted \(P(S)\)-tiling is the
standard \(\{6,3\}\) honeycomb.

\paragraph{Proof}

Lemma 5.1 makes the tiling a locally finite edge-to-edge trivalent
hexangulation.

Replace every physical face by a unit regular hexagon and glue corresponding
full sides. The topological realization of this abstract complex is
homeomorphic to the original polygon tiling and hence to \(\mathbb R^2\).
At every vertex, three sectors of angle \(2\pi/3\) meet. The regularized
surface is therefore locally Euclidean, including at its vertices.

The regularized path metric is complete. Triangulate every physical polygon
and its regular replacement by joining its barycenter to its six vertices.
On each prototile class, the resulting cellwise map and its inverse are
piecewise affine and Lipschitz. There are only finitely many prototile
classes, so one constant works for every triangle. The cellwise maps agree
on glued edges and give a global bi-Lipschitz homeomorphism from the physical
Euclidean plane to the regularized intrinsic surface. Completeness of the
physical plane therefore implies completeness of the regularized metric.

Choose an oriented frame in one face and analytically continue it across
edges. Since the surface is simply connected and flat, this gives a
developing local isometry
\[
  \operatorname{dev}:M\longrightarrow\mathbb R^2.
\]
A local isometry from a complete connected Euclidean surface to
\(\mathbb R^2\) has the path-lifting property for complete geodesics and is
a covering map. Because \(\mathbb R^2\) is simply connected,
\(\operatorname{dev}\) is a global Euclidean isometry. Its regular hexagonal
faces form the standard honeycomb. \(\square\)

The completeness step rules out a developing map whose image is only a
proper open subset of the plane. The covering step then supplies injectivity,
not merely local flatness.

\section*{6. Decoding every shape-only tiling}

\subsection*{Lemma 6.1. The line-graph star lemma}

At every vertex of an arbitrary \(P(S)\)-tiling, the three incident physical
sides recover three graph edges of \(G\) with one common endpoint state.

\paragraph{Proof}

Side length identifies a graph edge. Each pair among the three physical
sides occurs as a corner of one of the three incident polygon faces. The
corona that defined that polygon shows that the corresponding pair of graph
edges shares an endpoint.

The three graph edges are distinct. Otherwise an incident face would contain
two equal consecutive graph edges, contradicting the simplicity of its phase
face cycle.

Thus the three graph edges form a triangle in the line graph \(L(G)\). In a
simple graph, three pairwise incident edges either share one common vertex or
form a triangle. The graph \(G\) is bipartite, so it has no triangle. The
three edges form a star at one state \(q\). \(\square\)

\subsection*{Lemma 6.2. A physical edge recovers opposite state endpoints}

If a physical side decodes to \(\{q,r\}\in E(G)\), its two physical endpoints
decode to \(q\) and \(r\), in one of the two orders.

\paragraph{Proof}

Choose either incident polygon. In its corona cycle, the two ends of the side
are \(q\) and \(r\). At the physical endpoint corresponding to \(q\), the two
consecutive graph edges of that polygon share \(q\). Two distinct edges in a
simple graph have at most one common endpoint, so the common star state from
Lemma 6.1 is \(q\). The other endpoint similarly recovers \(r\). A reflected
placement may exchange the physical order, but cannot identify the two
states. \(\square\)

\subsection*{Lemma 6.3. The phase map is the universal honeycomb cover}

The decoded state labels define a cellular covering from the physical
honeycomb to \(H_2\).

\paragraph{Proof}

Map each physical vertex to the phase vertex contained in its decoded state,
each physical edge to the phase edge determined by its graph edge, and each
physical face to the phase face contained in its corona cycle.

At a physical vertex, the three incident graph edges project to distinct
phase edges. If two projected to the same phase edge, the physical face
between them would map two consecutive boundary edges to one edge of the
simple cycle \(C_{ij}\). Thus the vertex star maps bijectively to a phase
vertex star.

It remains to exclude folding along an edge. Suppose the two physical faces
adjacent to one edge mapped to the same phase face. At either endpoint, the
two other physical rays belonging to those faces would then map to the same
adjacent boundary edge of that phase face. This contradicts injectivity of
the vertex-star map. Hence the two sides of every physical edge map to the
two different incident phase faces.

The cellular map is now a local homeomorphism in face interiors, edge
interiors, and vertex stars.

For completeness, the image is all of \(H_2\). Starting from any image
vertex, bijectivity of its physical vertex star supplies all three incident
phase edges and their opposite endpoints. Iterating through the connected
one-skeleton reaches every phase vertex and edge. Along each lifted edge,
the no-folding argument supplies its two distinct incident phase faces, so
every face is reached as well.

Now take barycentric subdivisions of the physical complex and \(H_2\). A
vertex in either subdivision represents a face, an edge, or a vertex of the
original complex. Face recognition, the two-sided edge result, and the
bijective vertex-star result say exactly that the link of each physical
barycentric vertex maps isomorphically to the corresponding target link.
Consequently, the inverse image of every open barycentric star is a disjoint
union of open barycentric stars, each mapped homeomorphically onto the
target star. This is the cellular covering criterion. The phase map is
therefore a covering of \(H_2\). By Lemma 5.2 its domain is simply connected,
so it is the universal cover. \(\square\)

The local orientation sign of a covering between oriented surfaces is a
locally constant function with values in \(\{1,-1\}\). The physical
honeycomb is connected, so this sign is constant. Reflections are therefore
handled by one global choice rather than a hidden local chirality rule.

\subsection*{Lemma 6.4. Polygon placements have one global orthogonal frame}

There is one \(R\in O(2)\) such that the orthogonal part of every tile
placement is \(R\).

\paragraph{Proof}

A graph edge has one fixed canonical black-to-white vector because its state
endpoints include their phase vertices. The two canonical phase faces
incident to that edge traverse this vector in opposite directions and lie on
opposite sides of its line.

Consider two physical polygons sharing the corresponding side. After
orienting the canonical vector consistently, their orthogonal parts both map
it to the same physical vector. Two orthogonal maps agreeing on a nonzero
vector are either equal or differ by reflection in its line. In the second
case, the two canonical face interiors, initially on opposite sides, are
placed on the same physical side of the shared edge. Their interiors overlap
near the edge. This is forbidden in a tiling.

The orthogonal parts agree across every adjacency. The honeycomb adjacency
graph is connected, so one \(R\) applies to every tile. \(\square\)

\subsection*{Proposition 6.5. Source reconstruction}

Every unrestricted \(P(S)\)-tiling determines a point of \(S\), up to a
global deck-coordinate choice and the global frame \(R\).

\paragraph{Proof}

Lemma 6.3 identifies the physical cell complex with the universal cover of
\(H_2\). Use its deck coordinates to index physical faces by
\((i,j)\in\mathbb Z^2\). The six states around a physical face came from one
legal corona and therefore give the same central symbol. Lemma 6.2 makes
these state assignments agree across physical edges.

The phase edge specifies whether an adjacency is east-west, north-south, or
the unconstrained third direction. The legal triples and coronas impose
\(R_E\) and \(R_N\) on the first two directions. The reconstructed symbols
therefore define \(x\in S\). Every step used only side lengths, polygon
incidence, and the unmarked phase geometry encoded by those lengths.

Fix one physical vertex and its canonical lifted phase vertex. Lemma 6.4
gives one orthogonal part \(R\), and Lemma 6.2 gives the canonical vector
across every incident edge. Induction through the connected honeycomb now
shows that every physical vertex has the form
\[
  R\bigl(u+(z(q(u)),0)\bigr)+t
\]
for one translation \(t\). Thus the arbitrary tiling is exactly the
deformation associated with the reconstructed source point, up to \(R\) and
\(t\).
\(\square\)

Proposition 6.5 proves the soundness half of Theorem 1.1. It also classifies
every component of the shape-admitted tiling space. No selected generator
orbit is substituted for the complete space.

\section*{7. Constructing a tiling from every source point}

The converse requires a global deformation statement. A large finite patch
does not suffice.

\subsection*{Lemma 7.1. Quantitative honeycomb deformation}

Let \(x\in S\). Move each vertex of the affine honeycomb horizontally by the
displacement of its collared state. Add the barycenter of every face, move it
by the average of the six incident vertex displacements, and extend the map
affinely over the six center-fan triangles. The resulting map
\[
  F_x:\mathbb R^2\longrightarrow\mathbb R^2
\]
is a proper orientation-preserving local homeomorphism at bounded distance
from the identity.

\paragraph{Proof}

Every honeycomb vertex receives one state determined by the adjacent source
faces, so neighboring faces move a shared vertex to the same point. Every
vertex and barycenter displacement belongs to
\([0,\eta]\times\{0\}\).

For each center-fan triangle, direct expansion of its determinant gives
\[
  \det(F_x(v_k)-F_x(c),F_x(v_{k+1})-F_x(c))
  \geq 3-3\eta>0.
\]
The ray-star calculation following Lemma 4.2 gives the same lower bound for
every consecutive pair among the three deformed edge rays at a honeycomb
vertex. Thus every face fan stays embedded and positively oriented, while
every deformed three-face star keeps cyclic degree one. These facts make the
piecewise-affine extension a local homeomorphism in triangle interiors,
across triangulation edges, across tile edges, at face centers, and at
honeycomb vertices.

On every triangle, \(F_x-\operatorname{id}\) is a convex combination of its
vertex displacements. Hence
\[
  \lVert F_x(y)-y\rVert\leq\eta
\]
for every \(y\in\mathbb R^2\). If \(K\) is compact, then
\[
  F_x^{-1}(K)\subseteq
  \{y:\operatorname{dist}(y,K)\leq\eta\}.
\]
The preimage is closed and bounded, hence compact. Thus \(F_x\) is proper.
\(\square\)

\subsection*{Proposition 7.2. Completeness of the compiler}

If \(S\neq\varnothing\), then \(P(S)\) tiles \(\mathbb R^2\).

\paragraph{Proof}

Choose \(x\in S\). A proper local homeomorphism from the plane to the plane
has open image. Properness also makes its image closed. The image is
nonempty, so connectedness of \(\mathbb R^2\) makes it surjective. It is a
covering map, and the simply connected target has no nontrivial connected
cover. Hence \(F_x\) is a homeomorphism.

The images under \(F_x\) of the affine honeycomb faces cover the plane with
disjoint interiors. Each image face is exactly the polygon associated with
its legal corona. These images form a \(P(S)\)-tiling. \(\square\)

Combining Proposition 7.2 with Proposition 6.5 gives
\[
  P(S)\text{ tiles }\mathbb R^2
  \quad\Longleftrightarrow\quad
  S\neq\varnothing.
\]
This proves the tileability and decoding parts of Theorem 1.1.

\section*{8. Consequences}

\subsection*{8.1. Undecidable tileability}

Nonemptiness of nearest-neighbor \(\mathbb Z^2\)-shifts of finite type is
undecidable. The construction of \(P(S)\) uses finite enumeration and
rational arithmetic. Therefore an algorithm deciding unrestricted
tileability for finite sets of strictly convex rational all-obtuse hexagons
would decide whether \(S\) is nonempty. This proves Corollary 1.2.

The reduction permits rotations and reflections and does not assume
edge-to-edge contacts. The latter property is a conclusion of Lemma 5.1.
The decision-problem encoding admits the empty finite list, as specified in
Section 2.2.
After a common multiplication by \(2D\), every output coordinate is an
integer. The number of polygons is not bounded independently of the source.

\subsection*{8.2. Strong aperiodicity}

Let \(S\) be strongly aperiodic and suppose a \(P(S)\)-tiling
\(\mathcal T\) had translational period \(t\neq0\). The regularization and
decoding are local and commute with translations. Thus \(t\) preserves the
regular affine honeycomb obtained from \(\mathcal T\). It is a nonzero
honeycomb lattice vector and induces a nonzero element of
\(\mathbb Z^2\) in the deck coordinates of Lemma 6.3. Proposition 6.5 then
gives a point of \(S\) with that period, a contradiction. This proves
Corollary 1.3.

\subsection*{8.3. From ordinary aperiodicity to strong aperiodicity for a planar SFT}

For completeness, suppose a two-dimensional SFT has a configuration with
one nonzero period vector \(p\). Quotienting by \(p\) gives a nonempty
one-dimensional SFT over a finite alphabet of transverse strips. Every
nonempty one-dimensional SFT contains a periodic point, obtained from a
directed cycle in its finite presentation. Lifting that point gives a
doubly periodic configuration of the original SFT. Consequently, a Wang
shift with no doubly periodic configuration has no singly periodic
configuration either.

\section*{9. An explicit Kari-Culik instance}

\subsection*{9.1. The 13 source tiles}

Write a Wang tile as \((S,W,N,E)\). The six division tiles, with multiplier
\(1/3\), are

\begin{center}
\resizebox{\textwidth}{!}{%
\begin{tabular}{llll}
\toprule
\(S\) & \(W\) & \(N\) & \(E\) \\
\midrule
\(1\) & \(0/3\) & \(0\) & \(1/3\) \\
\(1\) & \(1/3\) & \(0\) & \(2/3\) \\
\(1\) & \(2/3\) & \(1\) & \(0/3\) \\
\(2\) & \(0/3\) & \(0\) & \(2/3\) \\
\(2\) & \(1/3\) & \(1\) & \(0/3\) \\
\(2\) & \(2/3\) & \(1\) & \(1/3\) \\
\bottomrule
\end{tabular}%
}
\end{center}

The seven doubling tiles, with multiplier \(2\), are

\begin{center}
\resizebox{\textwidth}{!}{%
\begin{tabular}{llll}
\toprule
\(S\) & \(W\) & \(N\) & \(E\) \\
\midrule
\(0\) & \(-1\) & \(0'\) & \(-1\) \\
\(0\) & \(0\) & \(0'\) & \(0\) \\
\(0\) & \(0\) & \(1\) & \(-1\) \\
\(1\) & \(-1\) & \(1\) & \(0\) \\
\(1\) & \(-1\) & \(2\) & \(-1\) \\
\(1\) & \(0\) & \(2\) & \(0\) \\
\(0'\) & \(0\) & \(1\) & \(-1\) \\
\bottomrule
\end{tabular}%
}
\end{center}

The colors \(0\) and \(0'\) are distinct. The symbols \(0/3,1/3,2/3\)
are formal colors.

\subsection*{9.2. An exact bi-infinite source point}

Set
\[
  \alpha_0=\sqrt2
\]
and define
\[
\alpha_{j+1}=
\begin{cases}
2\alpha_j,&\alpha_j<1,\\
\alpha_j/3,&\alpha_j>1.
\end{cases}
\]
The orbit stays in \([1/3,2)\). Its inverse is
\[
\alpha_{j-1}=
\begin{cases}
3\alpha_j,&\alpha_j<2/3,\\
\alpha_j/2,&\alpha_j>2/3.
\end{cases}
\]
Every \(\alpha_j\) is a nonzero rational multiple of \(\sqrt2\), so the orbit
never meets a rational branch boundary. The recurrence is therefore
unambiguous for every \(j\in\mathbb Z\).

For \(n,j\in\mathbb Z\), put
\[
  A_j(n)=\lfloor n\alpha_j\rfloor
\]
and let
\[
  \lambda_j=\frac{\alpha_{j+1}}{\alpha_j}\in\{2,1/3\}.
\]
Define
\[
\begin{aligned}
S_{n,j}&=A_j(n)-A_j(n-1),\\
N_{n,j}&=A_{j+1}(n)-A_{j+1}(n-1),\\
W_{n,j}&=\lambda_jA_j(n-1)-A_{j+1}(n-1),\\
E_{n,j}&=\lambda_jA_j(n)-A_{j+1}(n).
\end{aligned}
\]
On a division row, retain \(W\) and \(E\) in formal thirds notation. A zero
on the south side of row \(j\) is labeled \(0'\) exactly when
\[
  \alpha_{j-1}\in[1/3,1/2],
\]
and a zero on the north side is labeled \(0'\) exactly when
\[
  \alpha_j\in[1/3,1/2].
\]
All other numerical zeros are labeled \(0\).

\subsection*{Lemma 9.1. Exact floor classification}

For every \((n,j)\in\mathbb Z^2\), the four labels above form exactly one of
the 13 Kari-Culik tiles.

\paragraph{Proof}

Fix a row and write \(\alpha=\alpha_j\),
\(\beta=\alpha_{j+1}\).

First suppose \(\alpha<1\), so \(\beta=2\alpha\). For an integer \(k\), write
\[
  k\alpha=A+t,\qquad A\in\mathbb Z,\quad 0\leq t<1.
\]
Then
\[
  2A-\lfloor2k\alpha\rfloor=-\lfloor2t\rfloor\in\{-1,0\}.
\]
It follows that \(W,E\in\{-1,0\}\), while \(S\in\{0,1\}\) and
\(N\in\{0,1,2\}\). The definitions give the exact identity
\[
  2S+W=N+E.
\]
These ranges and this identity leave precisely the six numerical patterns
\[
\begin{gathered}
(0,-1,0,-1),\quad (0,0,0,0),\quad (0,0,1,-1),\\
(1,-1,1,0),\quad (1,-1,2,-1),\quad (1,0,2,0).
\end{gathered}
\]

If \(N=0\), then \(\beta<1\), hence
\(\alpha\in[1/3,1/2)\); the north zero is therefore \(0'\). If
\((S,W,N,E)=(0,0,1,-1)\), the south zero may occur on either side of
\(\alpha=2/3\). Indeed, the inverse branch is \(3\alpha\) below \(2/3\)
and \(\alpha/2\) above \(2/3\), so a south zero is \(0'\) exactly when
\(\alpha>2/3\). This gives the two distinct table entries
\((0,0,1,-1)\) and \((0',0,1,-1)\). The remaining south zeros in the two
\(N=0\) patterns are unprimed because there \(\alpha<1/2\). Thus the seven
doubling tiles are exactly the possible formal patterns.

Now suppose \(\alpha>1\), so \(\beta=\alpha/3\). For an integer \(k\), write
\[
  \lfloor k\alpha\rfloor=3\lfloor k\alpha/3\rfloor+r_k,
  \qquad r_k\in\{0,1,2\}.
\]
Then \(W=r_{n-1}/3\) and \(E=r_n/3\) as formal thirds. Also
\[
  S\in\{1,2\},\qquad N\in\{0,1\},
\]
and
\[
  \frac13S+W=N+E.
\]
Equivalently, if \(W=w/3\) and \(E=e/3\),
\[
  S+w=3N+e.
\]
The stated ranges leave exactly the six division rows in the table. A north
zero on a division row is unprimed because \(\alpha>1\). This proves
exhaustion and uniqueness. \(\square\)

\subsection*{Lemma 9.2. Matching and occurrence}

The floor construction is a valid bi-infinite Kari-Culik configuration. All
13 base tile types occur.

\paragraph{Proof}

Horizontal matching is an identity:
\[
  E_{n,j}=W_{n+1,j}.
\]
Vertical matching follows from
\[
  N_{n,j}=S_{n,j+1}.
\]
The prime convention is also identical on the shared vertical edge: both
sides use the condition
\(\alpha_j\in[1/3,1/2]\). Lemma 9.1 supplies one table tile at every lattice
site, including negative \(n\) and negative \(j\).

For the final assertion, write
\[
  \alpha_j=q_j\sqrt2
\]
and abbreviate
\[
  (a,b;c,d)=
  (A_j(n-1),A_j(n);A_{j+1}(n-1),A_{j+1}(n)).
\]
Direct exact evaluation gives one occurrence of every type:

\begin{center}
\resizebox{\textwidth}{!}{%
\begin{tabular}{lllll}
\toprule
Tile & \((n,j)\) & \(q_j\) & \((a,b;c,d)\) & Resulting \((S,W,N,E)\) \\
\midrule
1 & \((-1,0)\) & \(1\) & \((-3,-2;-1,-1)\) & \((1,0/3,0,1/3)\) \\
2 & \((2,0)\) & \(1\) & \((1,2;0,0)\) & \((1,1/3,0,2/3)\) \\
3 & \((-2,-3)\) & \(3/4\) & \((-4,-3;-2,-1)\) & \((1,2/3,1,0/3)\) \\
4 & \((-2,3)\) & \(4/3\) & \((-6,-4;-2,-2)\) & \((2,0/3,0,2/3)\) \\
5 & \((-2,0)\) & \(1\) & \((-5,-3;-2,-1)\) & \((2,1/3,1,0/3)\) \\
6 & \((3,0)\) & \(1\) & \((2,4;0,1)\) & \((2,2/3,1,1/3)\) \\
7 & \((-3,-2)\) & \(1/4\) & \((-2,-2;-3,-3)\) & \((0,-1,0',-1)\) \\
8 & \((4,-2)\) & \(1/4\) & \((1,1;2,2)\) & \((0,0,0',0)\) \\
9 & \((-1,1)\) & \(1/3\) & \((-1,-1;-2,-1)\) & \((0,0,1,-1)\) \\
10 & \((-1,-1)\) & \(1/2\) & \((-2,-1;-3,-2)\) & \((1,-1,1,0)\) \\
11 & \((-2,-1)\) & \(1/2\) & \((-3,-2;-5,-3)\) & \((1,-1,2,-1)\) \\
12 & \((-2,2)\) & \(2/3\) & \((-3,-2;-6,-4)\) & \((1,0,2,0)\) \\
13 & \((-3,-1)\) & \(1/2\) & \((-3,-3;-6,-5)\) & \((0',0,1,-1)\) \\
\bottomrule
\end{tabular}%
}
\end{center}

For \(\alpha_j=(p/q)\sqrt2\) and \(k>0\), every entry is certified by
\[
  A_j(k)=
  \left\lfloor
    \frac{\lfloor\sqrt{2k^2p^2}\rfloor}{q}
  \right\rfloor,
  \qquad
  A_j(-k)=-A_j(k)-1.
\]
The second identity uses irrationality of \(k\alpha_j\). The needed row
coefficients are
\[
q_{-3}=\frac34,\quad
q_{-2}=\frac14,\quad
q_{-1}=\frac12,\quad
q_0=1,\quad
q_1=\frac13,\quad
q_2=\frac23,\quad
q_3=\frac43.
\]
This finite table proves occurrence. It is not being used in place of the
symbolic all-site classification in Lemma 9.1. \(\square\)

Culik proves that this Wang set admits no periodic tiling. The
one-period-to-two-period argument in Section 8.3 makes its full Wang shift
strongly aperiodic. Corollary 1.3 therefore applies.

\subsection*{9.3. Size of the compiled instance}

For the phase-two Kari-Culik compilation, exact enumeration gives
\[
\begin{array}{c|r}
\text{quantity}&\text{value}\\ \hline
\text{source symbols}&13\\
\text{black vertex states }B&568\\
\text{white vertex states }W&600\\
\text{state-graph edges}&3632\\
\text{legal coronas}&4956\\
\text{congruence classes}&4956\\
D&1395916800.
\end{array}
\]
Thus the explicit shape-only protoset contains 4,956 pairwise noncongruent
strictly convex all-obtuse rational hexagons.

\section*{10. Exact artifact}

The retained artifact is

\begin{Verbatim}[fontsize=\small]
workbench/kc-convex-candidate.mjs
\end{Verbatim}

It uses \texttt{\detokenize{BigInt}} arithmetic for source floors, displacement coordinates,
cross products, dot products, squared lengths, and multiplier identities.
For a chosen winding modulus and phase modulus it enumerates every source
symbol, collared vertex state, graph edge, and legal corona. For the
unlifted Kari-Culik instance, replay:

\begin{Verbatim}[fontsize=\small]
node workbench/kc-convex-candidate.mjs 1 2 100
\end{Verbatim}

The principal output is:

\begin{Verbatim}[fontsize=\small]
alphabetSize                       13
blackCount                        568
whiteCount                        600
graphEdgeCount                   3632
coronaCount                      4956
congruenceClassCount             4956
exactConvexityChecks            29736
exactVectorConsistencyChecks     3632
exactUniqueLengthChecks          3632
coordinateFingerprintSha256
9be30b636e6a103c9360b5ffc7043e7d255df1a7e317da5369d86eb8c79e5fbd
\end{Verbatim}

The artifact checks:

\begin{enumerate}
\item the 13 multiplier identities in exact thirds units;
\item the exact bi-infinite source formulas on the requested finite audit
   window;
\item every face cycle is simple;
\item every exact corner cross product is positive;
\item every exact interior-ray dot product is negative;
\item every graph edge has one canonical black-to-white vector;
\item every graph edge has one squared length;
\item different graph edges have different squared lengths; and
\item distinct corona cycles remain distinct under the dihedral action.
\end{enumerate}

The SHA-256 value is a reproducibility fingerprint over the ordered
scaled-integer coordinate enumeration. It is not an independent proof
certificate. The symbolic arguments in Sections 4, 6, 7, and 9 prove the
infinite statements. The finite artifact checks the consequential finite
enumerations and catches implementation drift.

\section*{11. Effectivity and scope}

The compiler is a direct finite algorithm:

\begin{enumerate}
\item enumerate locally legal black and white state triples;
\item enumerate seven-face legal coronas in each of four phases;
\item construct the simple compatibility graph \(G\);
\item index its black and white vertices;
\item assign the mixed-radix rational displacements;
\item emit one exact hexagon per dihedral corona class.
\end{enumerate}

The output cardinality depends on the source alphabet and its local language.
This theorem supplies no fixed bound on the number of polygons. It is a
many-polygon compiler and does not settle tileability for any fixed number
of convex polygons.

If \(a=|A|\), then the four phases give the coarse bounds
\[
  B\leq4a^3,\qquad W\leq4a^3,\qquad |P(S)|\leq4a^7.
\]
These bounds are sufficient for effectivity and are not intended to be
sharp. When \(B=0\) or \(W=0\), the algorithm returns the empty protoset
before defining \(D\).

The construction is also unsuitable as a practical ceramic inventory in its
current form. Its edge differences may be far below ordinary manufacturing
tolerances, and the Kari-Culik instance has thousands of prototiles. Those
facts do not affect the exact tileability theorem, but they delimit its
material interpretation.

\section*{12. Relation to prior work}

Robinson's 1971 paper announces variable-cardinality undecidability for
convex lattice-coordinate pentagons or hexagons even with rotations and
reflections. It does not present that construction in the cited article.
This prior announcement prevents a broad priority claim for convex-polygon
undecidability.

Demaine and Langerman give a proved unrestricted-contact compiler using
three nonconvex polygons. Stade gives a two-polygon compiler with a highly
concave carrier. Yang and Zhang study translation-only tilings by
orthogonally convex polyominoes and ask fixed-cardinality questions. These
results have different output classes and cardinality goals. The result here
isolates a direct all-obtuse convex-hexagon mechanism.

Convex hexagon tilings can have irregular and non-edge-to-edge behavior in
other angle regimes. The present contact reduction uses the strict lower
bound \(\pi/2\) at every corner, rather than a general classification of
convex hexagon tilings.

\section*{13. Limitations and proof boundary}

The construction advances an output-class and forcing theorem. It does not
claim a new symbolic hull when instantiated on a known SFT. A faithful
nonconvex compiler applied to the same source may produce a locally
equivalent symbolic presentation.

The explicit Kari-Culik instance is large and visually close to an affine
honeycomb. Its role is to instantiate the theorem with exact data and a
strongly aperiodic source.

The novelty search behind this draft is finite. Robinson's unpublished
announcement is a reason to describe the result as an explicit all-obtuse
construction and proof, without an unconditional historical priority claim.

\section*{References}

\begin{enumerate}
\item Robert Berger, \emph{The Undecidability of the Domino Problem}, Memoirs of the
   American Mathematical Society 66, 1966.
\item Karel Culik II, "An aperiodic set of 13 Wang tiles,"
   \emph{Discrete Mathematics} 160 (1996), 245-251.
   \url{https://doi.org/10.1016/S0012-365X(96)00118-5}
\item Raphael M. Robinson, "Undecidability and nonperiodicity for tilings of the
   plane," \emph{Inventiones Mathematicae} 12 (1971), 177-209.
   \url{https://doi.org/10.1007/BF01418780}
\item Erik D. Demaine and Stefan Langerman, "Tiling with Three Polygons Is
   Undecidable," SoCG 2025, article 39.
   \url{https://doi.org/10.4230/LIPIcs.SoCG.2025.39}
\item Jack Stade, "Two Tiling is Undecidable," 2025.
   \url{https://arxiv.org/abs/2506.11628}
\item Chao Yang and Zhujun Zhang, "Undecidability of Translational Tiling with
   Orthogonally Convex Polyominoes," 2025.
   \url{https://arxiv.org/abs/2506.12726}
\item Dirk Frettlöh, Alexey Glazyrin, and Zsolt Lángi, "Hexagon Tilings of the
   Plane That Are Not Edge-to-Edge," 2019.
   \url{https://arxiv.org/abs/1911.12752}
\end{enumerate}

\end{document}
